\section{Introduksjon}
\label{sec:introduksjon}
\subsection{Bakgrunn}
Prosjektet omfatter utvikling av en mekanisk lekebil med svinghjulsbasert energilagring. Arbeidet høsten 2026 skal resultere i en fungerende fysisk prototype og en rapport som dokumenterer systematisk produktutvikling, 3D-design, produksjon og grunnleggende testing, med et omfang tilsvarende semesteroppgaven i MAST1101.

\subsection{Problemstilling}
\begin{quote}
Hvordan kan en mekanisk lekebil med svinghjulsbasert energilagring utvikles gjennom systematisk produktutvikling, 3D-design og prototyping for å oppnå størst mulig kjøredistanse ved ett standardisert opptrekk?
\end{quote}

\subsection{Målsetninger}
Målet er å utvikle og sammenligne praktiske løsninger innenfor prosjektets rammer. Gjennomsnittlig kjøredistanse er hovedmålet, supplert med repeterbarhet, driftssikkerhet og rettlinjet kjøring.
% Konkretiser resultatmålene, gjerne som en kort punktliste med målbare mål.

\begin{comment}
\subsection{Effektmål}
% Beskriv hva prosjektet skal oppnå utover selve prototypen, for eksempel læringsutbytte innen produktutvikling og et grunnlag for det videre arbeidet våren 2027.
\end{comment}

\subsection{Begrensninger}
Dette semesteret omfatter utvikling av bilen, derunder konseptutvikling og konseptvalg, CAD-modellering, produksjon, montering og grunnleggende eksperimentell testing. Rapporten utvikles samtidig med prototypen, slik at designvalg, endringer og forsøk blir dokumentert fortløpende.

Detaljert dynamisk og kinematisk modellering, validering av modellen og modellbasert optimalisering av drivverket avgrenses til MAST2007 våren 2027, der den samme bilen brukes videre. En enkel numerisk simulering brukes likevel i høst, som et verktøy for å sammenligne alternativer i iterasjonene (seksjon~\ref{sec:simulering-verktoy}). Den brukes til å rangere løsninger, ikke til å forutsi kjøredistansen. Høstens tester skal evaluere funksjon og ytelse og gi et sammenligningsgrunnlag for videre arbeid.
